\DocumentMetadata{
  pdfversion=2.0,
  pdfstandard=ua-2,
  lang=en-US,
  tagging=on,
  tagging-setup={math/setup=mathml-SE,tabsorder=structure}
}
\documentclass[11pt,letterpaper]{article}
\usepackage[margin=0.68in,includefoot]{geometry}
\usepackage{newcomputermodern}
\usepackage{amsmath,amscd,mathtools}
\usepackage{tikz,adjustbox}
\providecommand{\square}{\mdlgwhtsquare}
\usepackage[hidelinks]{hyperref}
\hypersetup{
  pdftitle={Combinatorics PhD Exam, May 2005},
  pdfauthor={University of Florida Department of Mathematics},
  pdfsubject={Graduate mathematics examination},
  pdfdisplaydoctitle=true,
  bookmarksopen=true
}
\setcounter{secnumdepth}{0}
\setlength{\parindent}{0pt}
\setlength{\parskip}{0.28em}
\setlength{\emergencystretch}{3em}
\setlength{\leftmargini}{2.15em}
\setlength{\leftmarginii}{2.65em}
\setlength{\leftmarginiii}{3em}
\pagestyle{empty}
\raggedbottom
\allowdisplaybreaks
\NewTaggingSocketPlug{block/list/label}{examproblem}{%
  \tagstructbegin{tag=Lbl}\tagstructend%
  \tagstructbegin{tag=\LBody}%
  \tagstructbegin{tag=\ProblemHeadingTag,title-o={\ProblemHeadingText}}%
  \tagmcbegin{tag=\ProblemHeadingTag,actualtext={\ProblemHeadingText}}%
  #2%
  \tagmcend%
  \tagstructend%
}
\newenvironment{examproblems}[2]{%
  \def\ProblemHeadingTag{#2}%
  \AssignTaggingSocketPlug{block/list/label}{examproblem}%
  \begin{enumerate}%
  \setcounter{enumi}{#1}%
  \renewcommand{\labelenumi}{\textbf{\theenumi.}}%
  \setlength{\topsep}{0.35em}%
  \setlength{\partopsep}{0pt}%
  \setlength{\itemsep}{0.48em}%
  \setlength{\parsep}{0.18em}%
}{\end{enumerate}}
\newenvironment{examsubparts}{%
  \AssignTaggingSocketPlug{block/list/label}{default}%
  \begin{enumerate}%
  \setlength{\topsep}{0.18em}%
  \setlength{\partopsep}{0pt}%
  \setlength{\itemsep}{0.16em}%
  \setlength{\parsep}{0.1em}%
}{\end{enumerate}}
\newenvironment{examinstructions}{%
  \par\begingroup\itshape
}{%
  \par\endgroup\vspace{0.18em}
}
\makeatletter
\renewcommand\subsection{\@startsection{subsection}{2}{0pt}%
  {-1.25ex plus -.25ex minus -.1ex}%
  {0.55ex plus .1ex}%
  {\normalfont\large\bfseries}}
\renewcommand{\@maketitle}{%
  \newpage\null\vskip -1em%
  {\footnotesize\bfseries
    \href{https://www.ufl.edu/}{University of Florida}, \href{https://math.ufl.edu/}{Department of Mathematics}\par}%
  \vskip -0.5em%
  \tagstructbegin{tag=H1,title={Combinatorics PhD Exam, May 2005}}%
  \tagpdfparaOff%
  \tagmcbegin{tag=H1}%
    {\LARGE\bfseries\href{https://gma.math.ufl.edu/exams/combinatorics-phd/}{Combinatorics PhD Exam}, May 2005\par}%
  \tagmcend%
  \tagpdfparaOn%
  \tagstructend%
  \vskip 0.25em%
  \tagmcbegin{artifact}%
    \hrule height 1pt%
  \tagmcend%
  \vskip 1em%
}
\makeatother
\title{Combinatorics PhD Exam, May 2005}
\author{}
\date{}
\begin{document}
\maketitle
\thispagestyle{empty}
\begin{examinstructions}
Show your work.
\end{examinstructions}
\begin{examproblems}{0}{H2}
\def\ProblemHeadingText{Problem 1.}
\item
Prove that, for any \(n+1\) numbers chosen from the set \(\{1,2,3,\ldots,2n\}\), one divides another.
\def\ProblemHeadingText{Problem 2.}
\item
State the Multiplier Theorem and use it to find two normalized \((11,5,2)\)-difference sets. (Normalized means that the sum of its elements is~\(0\).)
\def\ProblemHeadingText{Problem 3.}
\item
Show that a plane graph for which every face has an even number of edges must be bipartite.
\def\ProblemHeadingText{Problem 4.}
\item
Prove that for a binary linear code~\(C\), either all the codewords have even weight or exactly half the codewords have even weight. Show by example that this is not true for all ternary codes. Prove that if~\(C\) is a self-dual binary linear code, then all the codewords have even weight.
\def\ProblemHeadingText{Problem 5.}
\item
This problem has three parts.
\begin{examsubparts}
\item[\textnormal{(a)}]
Find the number of cycles of length \(2n\) in \(K_{n,n}\).
\item[\textnormal{(b)}]
Find the number of~\(6\)-cycles in \(K_{m,n}\).
\item[\textnormal{(c)}]
Find the number of~\(6\)-cycles in the~\(n\)-cube.
\end{examsubparts}
\def\ProblemHeadingText{Problem 6.}
\item
Find a recurrence for the number \(a_n\) of ways of going up~\(n\) stairs if we may take one or two steps at a time. Find the generating function for the sequence of \(a_n\) and an explicit formula for \(a_n\).
\def\ProblemHeadingText{Problem 7.}
\item
This problem has three parts.
\begin{examsubparts}
\item[\textnormal{(a)}]
Prove that for a \((v,k,\lambda)\)-BIBD we have \(\lambda(v-1)=r(k-1)\).
\item[\textnormal{(b)}]
Consider a symmetric \((v,k,\lambda)\)-BIBD and denote the order by \(n:=k-\lambda\). Using calculus to minimize, prove that \(v\geq 4n-1\).
\item[\textnormal{(c)}]
If \(v=4n-1\), show that \((v,k,\lambda)=(4n-1,2n-1,n-1)\) or \((4n-1,2n,n)\). Provide a family of symmetric designs with one of these sets of parameters.
\end{examsubparts}
\def\ProblemHeadingText{Problem 8.}
\item
If~\(\lambda\) is an eigenvalue of (the adjacency matrix~\(A\) of) a graph with~\(n\) vertices, \(m\) edges and maximum degree~\(\Delta\), then prove that
\begin{examsubparts}
\item[\textnormal{(a)}]
\(\lambda\leq\Delta\).
\item[\textnormal{(b)}]
\(\lambda\leq\sqrt{\frac{2m(n-1)}{n}}\).

\par
Hint on (b): use the trace of~\(A\), \(A^2\) and an appropriate inequality.
\end{examsubparts}
\end{examproblems}
\end{document}
